\documentclass{article}
\usepackage{arxiv}
\usepackage[utf8]{inputenc}
\usepackage[T1]{fontenc}
\usepackage{hyperref,url,microtype}
\hypersetup{colorlinks=true,urlcolor=blue,linkcolor=black}
\renewcommand{\shorttitle}{PSYCON: behavioral research synopsis}
\renewcommand{\headeright}{BEHA working draft}
\renewcommand{\undertitle}{Separate two-page research synopsis}
\title{PSYCON: Contextual communication behavior and evidence-linked psychologist observations}
\author{Arjun Vijay Prakash \And Saksham Yadav}
\date{October 13, 2026}
\begin{document}
\maketitle
\begin{abstract}
PSYCON turns recorded group discussions into speaker-specific communication measurements and evidence-linked reports for psychologist review. The implemented workflow combines normalized audio, diarization, transcription, timestamped exchanges and a versioned observation rubric that separates numeric ratings, insufficient opportunity and missing answers. This creates a reviewable connection between a communication observation and its supporting recording. The source-score pilot uses 423 psychologist-provided ratings from forty participant-session examples across nine recordings and five known connected groups. Compact grouped development improved agreement from 57.21\% to 62.19\% and reduced mean absolute error from 0.562 to 0.490 on the same 402 ratings. Training-mode agreement was 65.67\%, so reliable predictive superiority remains a question for the next study. A bounded 51-configuration audit tested feature families with training-only preprocessing and inner grouped selection, reproduced saved predictions and corrected a demonstrated temporal-ordering defect. Nine normalized recordings retained source audio duration within 15 milliseconds. The project delivered an executable measurement workflow, reproducible evaluation artifacts and reports linked to source evidence. These retrospective results support the prototype's research feasibility; class imbalance, unresolved identities and unreviewed reference labels require independent re-rating and validation before behavioral or coaching-benefit claims.
\end{abstract}

\section{Introduction and behavioral question}
The question is whether measurements and context predict reviewed A--T observations on held-out people/sources, and whether earlier history improves supported reports. The rubric concerns discussion tracking, contribution structure, turn-sharing, challenge, and delivery change. These are operational areas, not diagnostic factors. Reduced speech after challenge needs the event, earlier behavior, later response, and opportunity. Source context and measurement validity matter because a plausible model description cannot supply its own human validation.

\section{Innovation and objectives}
The contribution links participant, source time, exchange, reviewer, and revision in an auditable measurement workflow. H1 predicts lower held-out item error than the training-median baseline. H2 predicts better blinded grounding/usefulness without increased overclaiming. A history-only experiment must isolate history from the evidence/references also added by full PSYCON.

\newpage
\section{Methodology and controls}
The unit is a participant within a discussion. A reviewer confirms person/participant/speaker mapping. Nineteen author-supplied cross-recording links reduce 85 slots to at most 66 people and impose six provisional recording groups; further overlaps remain unknown. Linked recordings must share a partition. Recordings retain context, time, permission, and hashes. Audio processing conservatively attributes speech and preserves evidence; timing/marker proxies do not prove intent.

Psychologists supply scores, opportunity, confidence, recording support, and evidence, followed by independent review. Disagreement remains available before adjudication. Positive scores need timestamps; event-linked items need event/response windows, and Q--T need an earlier same-session baseline. Item S lacks predictors; visual-only T cannot train the audio model. Version 5.0 imports remain separate from v4.0; transcription and validity review are pending.

Frozen datasets group repeated people and related sources before splitting. Training-only preprocessing precedes per-item logistic classification; validation selects regularization, with test data separate. Fitting minimums are engineering gates, not a power calculation. Nine private signed sheets document voluntary participation. The authors confirm recording/training permissions, applicable parental permission and prior approval in the complete records; these details and recording linkage have not been independently checked. Analysis needs item error against the training-median baseline, weighted agreement, exclusions, abstentions, and grouped uncertainty. Reports need blinded ratings, randomized order, and a primary outcome. Language, topic, moderation, opportunity, and recording quality are confounds.

\section{Current results, limits, and next study}
The audiovisual comparison retained 57 PSYCON profiles versus 51 per Nemotron mode out of 85 slots, not unique people. The source-score pilot uses forty examples, nine recordings, five known groups and 423 ratings across G, I and P1--P12. Training agreement of 97.9\% is resubstitution. Original grouped development matched 230/402 ratings (57.21\%, MAE 0.562) with twenty-one abstentions. Compact selection reached 250/402 (62.19\%, 0.490), versus training median 259/402 (64.43\%, 0.460) and mode 264/402 (65.67\%, 0.453). The previous agreement-selection experiment reached 247/402 (61.44\%, 0.498). A new bounded inner-group-selected pooled procedure exactly matched mode predictions, reducing score-3 recall from 18.45\% to 4.85\%. A later 51-configuration feature-family autopsy yielded inner-selected 251/402 agreement (62.44\%, MAE 0.488), failing replacement gates. Chronological sorting fixed a demonstrated retry-latency defect without rewriting frozen features or claiming an accuracy gain. Compact remains the research reference; no candidate was promoted. All historical predictions refitted exactly. Other procedures score all 423 with no abstentions: compact 61.23\%, median 63.36\%, mode/new 64.54\%. Score 0 is unobserved; all-rating counts for 1--4 are 24,24,109,266. Twenty-one unresolved participant-session identities and absent independent review, opportunity, confidence and structured evidence prevent independent generalisation or a numerical ceiling. No final test supports 85\%. The report example has no eligible history or blinded usefulness result, so H2 remains open.

Next work must verify transcriptions, permission/mapping, and independent rater reliability, then freeze an independent sample. A history-only ablation should hold transcript, evidence, references, model, and decoding fixed. Poor agreement, no baseline improvement, or increased unsupported claims would challenge the hypotheses. Hardware/3D/KiCad work is paused future measurement research. Voluntary communication review is a potential use; diagnosis, disciplinary scoring, causation, and worldwide benefit are unsupported.

\section{Acknowledgements and selected references}
School psychologists supplied observation reference material; open-source authors supplied speech tools. Private identities and school/location identifiers are omitted. AI assisted this working draft; authors must independently rewrite applicable submission text and verify claims, references, contributions, consent, and approvals. No poster is included.

\small
Sacks, Schegloff and Jefferson (1974), turn-taking, \href{https://doi.org/10.2307/412243}{doi:10.2307/412243}. Cronbach and Meehl (1955), construct validity, \href{https://doi.org/10.1037/h0040957}{doi:10.1037/h0040957}. Cohen (1968), weighted agreement, \href{https://doi.org/10.1037/h0026256}{doi:10.1037/h0026256}. Full references and engineering sources accompany the separate paper. Study framing follows the \href{https://www.societyforscience.org/isef/grand-award/criteria/}{ISEF science judging criteria}; human-study requirements follow the \href{https://www.societyforscience.org/isef/international-rules/human-participants/}{Human Participants rules}.
\end{document}
